\chapter[Rare Regions and Learned Transports]{Rare Posterior Regions and Learned Transports}
\label{ch:rare-regions-transport}

A posterior region can contribute little to an average training loss and still
matter greatly for computation. A narrow valley between modes is one example:
its probability may be small, but a smooth transport must represent the rapid
change in density across it. In financial applications, a small-probability
region can also matter directly because the loss conditional on that region is
large. These are different reasons to study rare regions. Estimating a tail
probability, learning a density, and constructing useful Hamiltonian coordinates
require different evidence.

Let $\gamma(x)$ be an evaluable unnormalized posterior density on
$\mathbb R^d$, $Z=\int\gamma(x)\,dx\in(0,\infty)$, and
$\pi(x)=\gamma(x)/Z$. Write $A$ for an event and
$a=\Pr_\pi(X\in A)$. We assume that the densities and proposal probabilities
used below can be evaluated where their corrections require them. Knowing
$\gamma$ does not mean that its integral over $A$, its normalizer, or all of
its modes are known.

\section[Precision from ordinary samples]{What ordinary samples can and cannot resolve}

For independent $X_i\sim\pi$, the estimator
$\widehat a=N^{-1}\sum_i\mathbf1_A(X_i)$ has
\begin{equation}
 \mathbb E\widehat a=a,\qquad
 \operatorname{Var}(\widehat a)=\frac{a(1-a)}{N},\qquad
 \frac{\operatorname{SE}(\widehat a)}{a}
 =\sqrt{\frac{1-a}{Na}}.
 \label{eq:rare-iid-error}
\end{equation}
The variance follows by summing $N$ independent Bernoulli variances and
dividing by $N^2$. Achieving relative standard error $\delta$ requires
$N\geq(1-a)/(a\delta^2)$. This is a statement about standard error, not a
finite-sample confidence guarantee. The probability of observing no event is
$(1-a)^N$. For stationary correlated samples, the variance instead contains
the autocovariances of the event indicator. Effective sample size for a mean
coordinate cannot simply be substituted for effective sample size of a rare
indicator; a trace with almost no events provides little information about
the latter's dependence.

Consider the benchmark
\begin{equation}
 \pi(x_1,x_2)=\left[\tfrac13\phi(x_1+5)+\tfrac23\phi(x_1-5)\right]\phi(x_2),
 \label{eq:rare-mixture-target}
\end{equation}
where $\phi$ and $\Phi$ are the standard normal density and distribution
function. Symmetry gives
\begin{equation}
 \Pr(|X_1|<b)=\Phi(5+b)-\Phi(5-b).
 \label{eq:rare-mixture-mass}
\end{equation}
The broad valley $b=2$ has mass $0.001349898$, whereas the illustrative central
interval $b=0.1$ has mass $3.093565\times10^{-7}$. At a 10 percent relative
standard error, independent sampling requires about $74{,}000$ observations
for the former and $323$ million for the latter. The interval $0.1$ illustrates
the scaling; it is not an estimated width of a particular learned error.

In the recorded October 1 benchmark, one HMC run visited the broad valley
twice in $24{,}000$ retained draws, both times in one chain. Independent exact
sampling would average $32.4$ visits. This observation establishes inadequate
evidence for the valley probability, not an asymptotic bias theorem. The
previous validation allowance included a term proportional to $\sqrt a$;
for small $a$, that allowance can exceed the quantity being estimated. Relative
probability accuracy must be specified in units of $a$, with uncertainty
appropriate to the actual estimator.

\section[Density approximation and transport derivatives]{Density approximation does not bound transport derivatives}

For an invertible differentiable map $x=T(z)$, the latent posterior is
\begin{equation}
 \pi_z(z)=\pi(T(z))|\det J_T(z)|,\qquad
 r(z)=J_T(z)^\mathsf T\nabla\log\pi(T(z))
      +\nabla\log|\det J_T(z)|+z.
 \label{eq:rare-pullback-score}
\end{equation}
The residual $r$ vanishes for a standard Gaussian latent posterior. The unknown
$Z$ disappears on differentiation. Both terms involving the map are necessary;
a large individual term can be harmless if the terms cancel correctly.

The increasing exact transport of the first marginal of
\eqref{eq:rare-mixture-target} obeys $F_1(T_1(z))=\Phi(z)$. Differentiation
and then a logarithmic derivative give
\begin{align}
 T_1'(z)&=\frac{\phi(z)}{\pi_1(T_1(z))},\label{eq:rare-quantile-derivative}\\
 (\log T_1')'(z)&=-z-(\log\pi_1)'(T_1(z))T_1'(z).
\end{align}
Substitution in \eqref{eq:rare-pullback-score} proves exact cancellation. At
$T_1(z)=0$, $z\simeq-0.430727$ and $T_1'(z)\simeq244{,}565$. A large
derivative is therefore not itself an error. A learned map may fail because
it does not reproduce the cancellation. This example is a marginal
construction, not a lower bound for every multidimensional map.

Small KL divergence alone cannot exclude such derivative problems. To see
this, define $p_k(z)=Z_k^{-1}\phi(z)\exp\{k^{-1/2}\sin(kz)\}$. Since
$e^{-k^{-1/2}}\leq Z_k\leq e^{k^{-1/2}}$, the absolute log-density ratio
is at most $2k^{-1/2}$, so both KL directions tend to zero. Nevertheless,
\begin{equation}
 \mathbb E_\phi\left[(\partial_z\log p_k(Z)+Z)^2\right]
 =k\,\mathbb E_\phi[\cos^2(kZ)]
 =\frac{k}{2}(1+e^{-2k^2})\longrightarrow\infty.
 \label{eq:rare-kl-score-counterexample}
\end{equation}
The last equality uses the Gaussian characteristic function. This is a logical
counterexample, not a fitted description of the benchmark's error. Uniform or
derivative accuracy requires additional regularity, an appropriate objective,
or independent checks in the regions that matter.

Random base-space score checks can also miss an entire mode. In
\eqref{eq:rare-mixture-target}, take the simple map
$T(z)=(z_1+5,z_2)$, which generates just the right Gaussian component.
Writing $\phi_2$ for the standard bivariate Gaussian density gives the exact
identity
\begin{equation}
 \frac{\pi_z(z)}{\phi_2(z)}
 =\frac23+\frac13 e^{-10z_1-50},\qquad
 r_1(z)=-\frac{10e^{-10z_1-50}}{2+e^{-10z_1-50}},\quad r_2(z)=0.
 \label{eq:rare-missing-mode-score}
\end{equation}
At ordinary Gaussian base points the exponential term is tiny, so the score
residual is tiny, although one-third of the target mass lies near $z_1=-10$.
The sampled log-density offset
$\log\pi_z(z)+\|z\|^2/2$ is then close to
$-\log(2\pi)+\log(2/3)=-2.243342$, instead of the globally Gaussian value
$-\log(2\pi)=-1.837877$. This absolute-offset comparison uses the known
normalization of this benchmark. It cannot be applied unchanged when the
posterior normalizer is unknown. The example proves that small residuals on
ordinary base draws can coexist with a missed mode; it does not diagnose the
mode coverage of every learned map without further evidence.

\section[Umbrella sampling and relative probabilities]{Umbrella sampling and the recovery of relative probabilities}

Umbrella sampling allocates separate simulations to overlapping regions. For
nonnegative bias functions $\psi_k$ with $S(x)=\sum_k\psi_k(x)>0$ on the
target support, define
\begin{equation}
 c_k=\mathbb E_\pi\psi_k(X),\qquad
 \pi_k(x)=\frac{\psi_k(x)\pi(x)}{c_k}.
 \label{eq:rare-umbrella-density}
\end{equation}
A common choice is
$\psi_k(x)=\exp\{-\kappa_k(\xi(x)-b_k)^2/2\}$, where $\xi$ identifies a
barrier or rare region. A simulation centered in a valley then observes that
valley frequently under $\pi_k$, even when it is rare under $\pi$.

The biased distribution can be simulated without knowing $c_k$. For example,
put $\ell_k=\log\gamma+\log\psi_k$ and propose the Langevin step
\begin{equation}
 Y=X+h\nabla\ell_k(X)+\sqrt{2h}\,\varepsilon,
 \qquad \varepsilon\sim N(0,I),\quad h>0.
 \label{eq:rare-window-mala}
\end{equation}
Its Gaussian proposal density is
$Q_k(y\mid x)=(4\pi h)^{-d/2}
\exp\{-\|y-x-h\nabla\ell_k(x)\|^2/(4h)\}$.
Accept with probability
$\min\{1,e^{\ell_k(y)-\ell_k(x)}Q_k(x\mid y)/Q_k(y\mid x)\}$.
The unknown normalizers cancel, and the accepted flux is symmetric after
multiplication by $\pi_k(x)Q_k(y\mid x)$. Omitting the proposal ratio would
generally change the invariant distribution. A rejection leaves $X$
unchanged. This correction establishes invariance; the simulation must still
explore its biased distribution adequately.

The unknown constants $c_k$ prevent simple pooling. The eigenvector method
for umbrella sampling (EMUS) resolves their ratios from overlapping simulations
\citep[sections II--IV]{thiede2016emus}. Define
\begin{equation}
 F_{ij}=\mathbb E_{\pi_i}\!\left[\frac{\psi_j(X)}{S(X)}\right],
 \qquad z_i=\frac{c_i}{\sum_k c_k}.
 \label{eq:rare-emus-overlap}
\end{equation}
Every row of $F$ sums to one. Moreover,
\begin{align}
 \sum_i c_iF_{ij}
 &=\sum_i\int\frac{\psi_i(x)\psi_j(x)}{S(x)}\pi(x)\,dx\\
 &=\int\psi_j(x)\pi(x)\,dx=c_j.
\end{align}
Hence $z^\mathsf TF=z^\mathsf T$ and $\sum_i z_i=1$. Connected overlap,
expressed by irreducibility of $F$, supplies a unique normalized solution.
For an integrable observable $g$,
\begin{equation}
 \mathbb E_\pi g(X)=
 \frac{\sum_i z_i\mathbb E_{\pi_i}[g(X)/S(X)]}
      {\sum_i z_i\mathbb E_{\pi_i}[1/S(X)]}.
 \label{eq:rare-emus-average}
\end{equation}
Indeed, multiplying each expectation by $c_i$ cancels its normalizer, and
summing cancels $S$. The denominator removes the common scale lost when
replacing $c$ by $z$. Setting $g=\mathbf1_A$ estimates the event probability;
setting $g=-\log q_\theta$ supplies a corrected training objective.

EMUS is related to multistate Bennett acceptance ratio estimation (MBAR),
as explained in the paper's section IV. Neither method makes disconnected or
unequilibrated windows informative. Section VII derives a central limit
theorem assuming within-window central limit theorems, independent windows,
positive limiting sample allocations, and irreducible overlap. Its asymptotic
variance has the form
\begin{equation}
 \sigma^2=\sum_i\frac{\chi_i^2}{\alpha_i},\qquad
 \alpha_i=N_i/\sum_jN_j.
 \label{eq:rare-emus-allocation}
\end{equation}
Here $\chi_i^2$ combines the sensitivity of the reported observable to window
$i$ with that window's integrated autocovariance. When all $\chi_i>0$,
minimizing this expression subject to $\sum_i\alpha_i=1$ gives
$\alpha_i\propto\chi_i$ by a Lagrange multiplier. A zero contribution gives
a boundary optimum; removing that window still requires preserving the
coverage and overlap assumptions. Thus useful allocation depends on the observable, not just on the
posterior mass of the window. These are asymptotic statements; short traces
cannot certify their own equilibration.

\section[Importance sampling for an observable]{Importance sampling directed at an observable}

Let $q$ be a normalized, samplable proposal with support covering $\pi$.
For $X_i\sim q$ and $w_i=\gamma(X_i)/q(X_i)$, the self-normalized estimate is
\begin{equation}
 \widehat\mu_g=\frac{\sum_iw_ig(X_i)}{\sum_iw_i}
 \longrightarrow\frac{\int\gamma(x)g(x)\,dx}{\int\gamma(x)\,dx}
 =\mathbb E_\pi g(X).
 \label{eq:rare-snis}
\end{equation}
The convergence follows from laws of large numbers for numerator and
denominator, under the corresponding integrability conditions. The ratio is
generally biased at finite $N$. Its asymptotic variance, when finite, is
\begin{equation}
 \sigma_g^2(q)=\int\frac{\pi(x)^2(g(x)-\mu_g)^2}{q(x)}\,dx.
 \label{eq:rare-snis-variance}
\end{equation}
To obtain this expression, subtract $\mu_g$ from
\eqref{eq:rare-snis} and multiply by $\sqrt N$. The numerator becomes
$N^{-1/2}\sum_iw_i(g(X_i)-\mu_g)$, whose expectation is zero. The
denominator $N^{-1}\sum_iw_i$ converges to $Z$. A central limit theorem
for the numerator, followed by this denominator limit, gives variance
$Z^{-2}\mathbb E_q[w^2(g-\mu_g)^2]$, which is
\eqref{eq:rare-snis-variance}. A finite value of this integral is an
additional assumption, not a consequence of support coverage alone.
Applying Cauchy--Schwarz shows that the oracle minimizing proposal is
$q^*(x)\propto\pi(x)|g(x)-\mu_g|$. For an event indicator with $0<a<1$, the integrals
of this unnormalized proposal inside and outside $A$ are both $a(1-a)$.
The optimal proposal therefore assigns half its probability to each side,
however small $a$ is. This derivation, discussed in Owen's chapter 9,
explains why matching the entire posterior is not the same problem as
estimating one rare probability \citep[section 9.2]{owen2013mc}.

The oracle proposal requires information that is ordinarily unknown. A
practical mixture can combine approximations around modes, proposals near the
event, and a full-support component. Every sample must be weighted by the
density of the complete mixture. Full support prevents one particular bias,
but does not establish finite variance or useful finite-sample overlap.
Proposal adaptation should use pilot draws and be frozen before an independent
assessment. If old and new proposal samples are combined, their distinct
sampling laws must enter an appropriate multiple-importance correction.

\section{Multilevel splitting and conditional exploration}

Choose nested events $A_0=\mathbb R^d\supset A_1\supset\cdots\supset A_m=A$.
The multiplication rule for conditional probabilities yields
\begin{equation}
 a=\prod_{j=1}^m a_j,\qquad
 a_j=\Pr_\pi(X\in A_j\mid X\in A_{j-1}).
 \label{eq:rare-splitting-product}
\end{equation}
Splitting estimates each factor with a population directed toward the next
event. Survivors are resampled and moved using a kernel invariant for
$\pi(\cdot\mid A_j)$. Adaptive multilevel splitting (AMS) chooses levels
from particle quantiles. The idealized AMS analysis of \citet{brehier2015ams} assumes exact
independent conditional draws; section 2.2 explicitly distinguishes that
assumption from approximate finite Metropolis mutation in a static problem.
Consequently, its unbiasedness and cost results must not be transferred to
an arbitrary approximate implementation.

The same populations can supply a weighted training bank. Let
$B_j=A_{j-1}\setminus A_j$ and $P_j=\Pr_\pi(A_j)$ with $P_0=1$. Then
\begin{equation}
 \mathbb E_\pi g
 =\sum_{j=1}^m(P_{j-1}-P_j)\mathbb E_\pi[g\mid B_j]
   +P_m\mathbb E_\pi[g\mid A_m].
 \label{eq:rare-splitting-shells}
\end{equation}
This is the law of total expectation over a disjoint partition. If stage
$j$ starts with $N$ equal-weight particles approximating
$\pi(\cdot\mid A_{j-1})$, each rejected shell particle carries weight
$\widehat P_{j-1}/N$; the surviving fraction defines
$\widehat P_j$. After the final mutation, terminal particles carry
$\widehat P_m/N$. The estimated shell masses telescope to one. Conservation
of this mass verifies bookkeeping, not correctness of the approximate
conditional samples. Resampling ancestry, mutation movement, repeated
populations, and uncertainty in the product remain necessary diagnostics.

\section[Local maps and corrected global moves]{Local learned maps and global corrected moves}

A single smooth map from a unimodal base may have difficulty representing a
narrow connection between distant modes. Gabri\'e, Rotskoff and
Vanden-Eijnden discuss mixtures of maps trained in distinct modes
\citep[section IV.E]{gabrie2022adaptiveflows}. If $q_k$ is the density induced
by map $T_k$ and $\eta_k>0$ with $\sum_k\eta_k=1$, form
\begin{equation}
 q_\eta(x)=\sum_k\eta_kq_k(x).
 \label{eq:rare-local-map-mixture}
\end{equation}
Draw a component and then a point from its map. Given current state $x$, a
global proposal $y\sim q_\eta$ is accepted with probability
\begin{equation}
 \alpha(x,y)=\min\left\{1,
 \frac{\gamma(y)q_\eta(x)}{\gamma(x)q_\eta(y)}\right\}.
 \label{eq:rare-mixture-mh}
\end{equation}
Detailed balance follows because the accepted probability flux,
\[
 \min\{\pi(x)q_\eta(y),\pi(y)q_\eta(x)\},
\]
is symmetric in $x,y$.
Composition with a target-invariant local kernel preserves invariance. In
particular, global proposals can cross from one mode to another without
following a continuous path through the valley.

Local training need not know the modes' posterior masses. Those masses are
recovered through corrected sampling, rather than set equal to the initial
numbers of walkers. Mixture weights may be adapted during warm-up by fitting
the current sampled distribution, then frozen before retained estimation.
For example, write $\eta_k=e^{a_k}/\sum_j e^{a_j}$ and let
$R_k(x)=\eta_kq_k(x)/q_\eta(x)$ be the component responsibility.
Differentiation gives
\begin{equation}
 \frac{\partial[-\log q_\eta(x)]}{\partial a_k}
 =\eta_k-R_k(x).
 \label{eq:rare-mixture-logit-gradient}
\end{equation}
A stochastic gradient step on the average of this expression adjusts the
mixture weights using sampled occupancy. If a fixed defensive component
has mass $d$, write $q=(1-d)\sum_k\eta_kq_k+dq_D$ instead. Defining
$R_k=(1-d)\eta_kq_k/q$ changes the gradient to
$\eta_k\sum_jR_j-R_k$; the defensive mass itself stays fixed. These are
cross-entropy derivatives, not independent estimates of unknown mode masses.
Frozen-kernel invariance does not prove convergence of an arbitrary indefinitely
adaptive procedure. The authors' section IV.C and Appendix E also state that
unknown modes should not be expected to be discovered automatically. Their
Appendix G.1 uses a mixture with separation ten and weights one-third and
two-thirds, and discusses the residual connection learned by a unimodal-base
flow. Successful mode mixing still does not imply a precise probability for
an extremely rare region between the modes.

\section[Deliberate exploration and NeuTra training]{Combining deliberate exploration with NeuTra training}

For a disjoint partition $B_k$ with masses $b_k$, forward-KL training minimizes
the cross entropy
\begin{equation}
 L(\theta)=-\mathbb E_\pi\log q_\theta(X)
 =-\sum_kb_k\mathbb E_{\pi(\cdot\mid B_k)}\log q_\theta(X).
 \label{eq:rare-stratified-training}
\end{equation}
Drawing $n_k$ points within each stratum gives the corrected estimator
\begin{equation}
 \widehat L(\theta)=-\sum_k\frac{b_k}{n_k}
     \sum_{i=1}^{n_k}\log q_\theta(X_{ki}).
 \label{eq:rare-stratified-estimator}
\end{equation}
Thus the training batch can deliberately contain many valley points without
pretending that the valley has large posterior mass. For an existing weighted
bank, sample within each stratum proportionally to its conditional weights
and give each selected point weight $b_k/n_k$. The total batch weight is then
exactly one; ordinary posterior resampling would undo the deliberate allocation.
Estimated stratum masses and imperfect conditional samples introduce their
own error. Overlapping umbrellas require \eqref{eq:rare-emus-average}, not
unweighted pooling or the disjoint formula applied without correction.

This allocation repairs the information available to the optimizer. It does
not change the small coefficient $b_k$ in the original objective. Increasing
that coefficient or adding a score penalty defines a different optimization
problem and must be stated explicitly. Reverse-KL refinement can subsequently
use the exact transformed target, but a decreasing objective still cannot
replace derivative and downstream checks.

More explicitly, for standard Gaussian base density $\phi_d$ and invertible
$T_\theta$, change of variables gives
$q_\theta(T_\theta(z))=\phi_d(z)/|\det J_{T_\theta}(z)|$. Substitution
in the reverse KL objective yields the computable loss
\begin{equation}
 L_{\mathrm{RKL}}(\theta)=\mathbb E_{z\sim\phi_d}\!\left[
  \log\phi_d(z)-\log\gamma(T_\theta(z))
  -\log|\det J_{T_\theta}(z)|\right].
 \label{eq:rare-rkl-refinement}
\end{equation}
It differs from $\mathrm{KL}(q_\theta\|\pi)$ by the constant
$-\log Z$, so the unknown posterior normalizer does not affect its gradient.
The weighted forward fit uses deliberately collected physical-space points;
this refinement instead draws Gaussian base points and maps them with the
current parameters. Consequently,
it does not keep presenting a rare physical region merely because that region
was represented in the earlier bank. Retaining forward checkpoints and checking
mode and valley behavior after refinement are therefore necessary to distinguish
a useful refinement from a loss decrease accompanied by lost coverage.

For the separable benchmark, the exact marginal CDF provides a direct
reference construction. Directed physical-space probes can inspect the valley
and then be mapped back to latent coordinates. Such probes complement random
Gaussian probes; neither their finite maximum nor a sample quantile bounds
the score everywhere. In higher dimensions, an informative coordinate,
overlap between windows, well-mixing conditional kernels, and tractable local
structure become substantive modeling and computational questions. These
methods allocate effort more deliberately than passive sampling, but none
guarantees discovery of every arbitrary rare feature from finitely many
density evaluations.

\section[Two-dimensional sampling evidence]{Evidence from the two-dimensional benchmark}

An October 2 experiment applied the four constructions to
\eqref{eq:rare-mixture-target} and to the transformed target
$(X_1,X_2+0.1(X_1^2-26))$. This transformation has unit Jacobian and
preserves the first-coordinate event probabilities. Mode locations and the
valley coordinate were supplied; mode probabilities and exact reference draws
were withheld from proposal training. Eight independent estimator runs per
method supplied the means and between-run standard errors in
Table~\ref{tab:rare-region-estimation}. Fixed-level splitting used finite
Metropolis mutation, so the ideal independent-conditional AMS theorem does
not apply to these estimates.

\begin{table}[htbp]
 \centering\small
 \caption{Estimated probabilities, with between-run standard errors in
 parentheses. The narrow event has exact probability
 $3.093565\times10^{-7}$. Local maps report the longer repair run;
 ``no visits'' supplies no usable estimate of relative precision.}
 \label{tab:rare-region-estimation}
 \begin{tabular}{llrrr}
 \hline
 Target & Method & $10^3\Pr(|X_1|<2)$ & $10^7\Pr(|X_1|<0.1)$ & $\Pr(X_1>0)$\\
 \hline
 Unwarped & Umbrella & 1.070 (0.162) & 2.469 (0.552) & 0.638 (0.065)\\
 & Importance & 1.296 (0.033) & 3.012 (0.034) & 0.670 (0.004)\\
 & Splitting & 1.341 (0.049) & 3.057 (0.138) & 0.665 (0.002)\\
 & Local maps & 1.373 (0.159) & no visits & 0.666 (0.003)\\
 Warped & Umbrella & 1.027 (0.131) & 2.343 (0.460) & 0.589 (0.067)\\
 & Importance & 1.298 (0.032) & 3.010 (0.033) & 0.669 (0.004)\\
 & Splitting & 1.343 (0.029) & 2.963 (0.130) & 0.666 (0.003)\\
 & Local maps & 1.694 (0.169) & no visits & 0.666 (0.003)\\
 \hline
 Exact & & 1.349898 & 3.093565 & 0.666667\\
 \hline
 \end{tabular}
\end{table}

Umbrella sampling, importance sampling, and splitting passed the declared
feasibility screen on both targets: every replication observed each event,
the standard error of each aggregate probability estimate was at most
20 percent of its exact probability, and its discrepancy was no greater than
the larger of four standard errors and 10 percent of that probability.
First and second moments were checked separately. This is a working
feasibility screen, not a simultaneous confidence guarantee or a statistically
established ranking. The sizeable uncertainty of the umbrella estimates is
visible in the table.

The local-map mixture passed the mode and moment checks but never visited
the narrow event, including after its longer warm-up and doubled retained
sample count. The latter run had 8192 retained states per replication. Even
65536 independent posterior draws would have only about 0.020 expected
narrow-event visits; independence is not assumed for these chains. The
absence of visits therefore fails the requested rare-probability measurement
without contradicting invariance of the corrected sampler. Exact iid
sampling with 4096 points per replication also missed this event, and the
unaugmented local Laplace importance proposal failed the precision screen.
These controls expose the cost of passive sampling and the value of directing
proposal mass toward the specified event.

For subsequent global-map training, the first replication was fixed in
advance as the weighted bank. Aggregate agreement over eight replications
does not certify this bank individually: the two umbrella banks, for example,
assigned the right mode probabilities 0.759 and 0.762. Stratified minibatches
preserved each bank's weights, including rare rows. They did not remove its
sampling error or replace the need to check the learned map and its HMC draws.

The subsequent test trained a global inverse autoregressive flow (IAF) from
each bank, using two training seeds per target and teacher: sixteen fits in
total. Each used a crossed width/learning-rate pilot, 1024 forward-KL updates,
and reverse-KL checkpoints at 256 and 1024 updates. Forty-eight checkpoints
produced finite 1000-point Gaussian-base diagnostics and directed valley
checks. Final median random-point residuals ranged from 0.052 to 0.574,
whereas the maxima on the directed valley line ranged from 10.252 to 380.096.
These describe different probe measures, not uniform error bounds. Several
sampled log-density offsets lay near the one-mode value in
\eqref{eq:rare-missing-mode-score}, which supplies an alternative explanation
for their small local residuals.

All twelve fits whose teachers passed the probability screen then underwent
fresh fixed-transport HMC tuning and sequential posterior assessment with
identity mass in latent coordinates. The earlier checkpoints remained
available. The assessment covered moments, mode occupancy, marginal CDF cuts,
and the broad valley; it did not demand precise ordinary-sampling estimation
of the $3.09\times10^{-7}$ central event. No fit passed within the declared
limits. Of 41 verified kernel members assessed, four reached retained
sampling and exhausted the 10000-draw-per-chain limit without the required
ESS and precision. Fourteen members failed warm-up because a chain stopped
moving; other failures involved convergence, missing event observations,
or an exhausted search allocation. Final confirmation data remained unopened.

The experiment therefore supports deliberate estimation of the specified
rare probabilities on these two targets, while leaving the global NeuTra
training problem unresolved. It does not establish a ranking among the three
viable estimators or an impossibility result for IAFs. The limited training
ladder, individual bank error, and bounded HMC search remain alternative
failure explanations. The exact marginal transport available here provides
a useful next control for separating a learned-coordinate failure from a
restriction in downstream qualification.
